\documentclass{article}
\usepackage{amsmath,amssymb,amsfonts}
\usepackage{graphicx}
\usepackage{amsmath,amssymb,amsfonts}
\usepackage{amsmath,amssymb,amsthm,physics}
\usepackage{graphicx}
\usepackage{parskip}
\usepackage{amsmath}
\usepackage{amssymb}
\usepackage{amsfonts}
\usepackage{mathtools}
\usepackage{physics}
\usepackage{amsthm}  % This provides theorem-like environments

\begin{document}
\title{Formal Mathematical Proof of the Holographic Galactic Coordinate System}
\author{Bryce Weiner}
\date{\today}

\section*{Formal Mathematical Proof of the Holographic Galactic Coordinate System}

\subsection*{Definitions and Fundamental Constants}

Let the following be the fundamental constants and definitions in the holographic galactic coordinate system:

\begin{align}
\gamma &\approx 1.89 \times 10^{-29} \text{ s}^{-1} && \text{Information generation rate} \\
\pi &\approx 3.14159265359 && \text{Mathematical constant} \\
\alpha &= \frac{1}{4\pi \times 256} && \text{Continuum bifurcation coupling constant}
\end{align}

\subsection*{Establishing the Reference Frame}

Define the reference frame as follows:

\begin{itemize}
    \item Origin: Location of the Great Cold Spot (GCS) in the cosmic microwave background
    \item Orientation: Aligned with the cosmic microwave background dipole anisotropy
\end{itemize}

Let the coordinates of the GCS be:
\begin{align}
l_\text{GCS} &= 209^\circ \\
b_\text{GCS} &= -57^\circ \\
z_\text{GCS} &= 0.0155
\end{align}

\subsection*{Constructing the Distance Scale}

Define the luminosity distance to an object as:
\begin{equation}
d_L = \frac{L_\text{standard}}{4\pi F_\text{observed}}
\end{equation}
where $L_\text{standard}$ is the standardized peak luminosity of Type Ia supernovae.

The distance from baryon acoustic oscillations (BAO) is given by:
\begin{equation}
d_\text{BAO} = \frac{r_d^\text{HU}}{
\theta_\text{BAO}(z)}
\end{equation}
where $r_d^\text{HU} = r_d^\text{ΛCDM} \left(1 - \gamma t_\text{drag}\right)$ is the modified BAO scale.

\subsection*{Mapping Galactic Structures}

Compile a catalog of Type Ia supernovae and BAO features with their positions and redshifts.

Convert the supernova positions $(l, b, z)$ to Cartesian coordinates $(X, Y, Z)$ using:
\begin{align}
X &= d_L \cos(b) \cos(l) \\
Y &= d_L \cos(b) \sin(l) \\
Z &= d_L \sin(b)
\end{align}

Similarly, convert the BAO positions to Cartesian coordinates using the BAO distance.

\subsection*{Incorporating Holographic Principles}

Define the velocity ratio between two elements as:
\begin{equation}
\frac{v_1}{v_2} = \sqrt{1 + \gamma t_1} / \sqrt{1 + \gamma t_2}
\end{equation}

The evolution of the R(Si II) ratio follows:
\begin{equation}
R(\text{Si II})(t) = R(\text{Si II})_0 \exp(-\gamma t)
\end{equation}

The locations of the node points $N_p(n)$ are given by:
\begin{equation}
N_p(n) = N_0 \exp(2\pi i /n)
\end{equation}
where $n$ is an integer.

\subsection*{Validating and Refining the Coordinate System}

Verify the consistency of the coordinate system by checking:
\begin{align}
\Delta T_\text{GCS} &= -\gamma V_\text{GCS} / c^3 \\
\frac{v_1}{v_2} &= \sqrt{1 + \gamma t_1} / \sqrt{1 + \gamma t_2} \\
R(\text{Si II})(t) &= R(\text{Si II})_0 \exp(-\gamma t) \\
N_p(n) &= N_0 \exp(2\pi i /n)
\end{align}

Update the coordinate system by incorporating new observational data, recomputing distances, and validating the consistency.

\end{document}